\begin{abstract}
Class-conditional conformal prediction calibrates a separate score threshold for every class
and therefore requires labelled calibration examples of every class. In long-tailed
deployments these examples are precisely what is missing, as most rare classes are consumed
by training and are left with few or no held-out examples. Existing class-conditional methods
are then either undefined for such a class or fall back to the shared marginal threshold. We
ask whether the per-class correction can instead be predicted from descriptors of the class
that require no labels at all.
%
We propose Predicted Class Correction (PCC), which regresses the offset between the per-class
and the marginal conformal quantile on geometric descriptors of the classifier head, shrinks
the prediction by a scalar selected on the labelled classes, and refits a single marginal
offset. Because the descriptors are functions of the model's weights, a class with zero
calibration examples still receives a threshold specific to it, and because only one scalar is
calibrated, the marginal coverage guarantee is retained
(Prop.~\ref{prop:marginal}).
%
On ImageNet, PCC improves worst-class coverage on classes with zero calibration examples by
$0.075$ at matched average set size, which is $70\%$ of an oracle that observes the test
labels, while the strongest of the seven competitors we evaluate gains $0.027$ and two are
undefined on all $300$ such classes. The improvement persists when $90\%$ of classes are held
out, leaving $100$ labelled classes out of $1000$.
%
We also study where the approach stops working, and find that the limit is a property of
the evaluation set. Below roughly $35$ test examples per class the oracle ceiling itself
falls to zero, so no method can be shown to help there. Restricting
evaluation to classes above that threshold recovers a significant improvement on
Pl@ntNet-300K, where the effect had appeared to be absent.
\end{abstract}
